% ========================================================
% BÀI TẬP TOÁN 5 - KẾT NỐI TRI THỨC VỚI CUỘC SỐNG
% Bài 28: Thực hành và trải nghiệm đo, vẽ, lắp ghép tạo hình
% Định dạng: BTTeX Part III - Trắc nghiệm điền số ngắn
% ========================================================

% Câu 1
\begin{ex}
  Ghép hai hình tam giác vuông cân bằng nhau có cạnh góc vuông $6\text{ cm}$ thành một hình vuông. Diện tích hình vuông đó là bao nhiêu xăng-ti-mét vuông?
  \shortans{36}
  \loigiai{
    Diện tích hình vuông là $6 \times 6 = 36\text{ cm}^2$.
  }
\end{ex}

% Câu 2
\begin{ex}
  Ghép hai hình tam giác đều có cạnh $10\text{ cm}$ lại với nhau để được một hình thoi. Chu vi của hình thoi đó là bao nhiêu xăng-ti-mét?
  \shortans{40}
  \loigiai{
    Chu vi hình thoi là $10 \times 4 = 40\text{ cm}$.
  }
\end{ex}

% Câu 3
\begin{ex}
  Một hình chữ nhật có kích thước $8\text{ cm} \times 4\text{ cm}$. Người ta cắt dọc theo đường nối trung điểm hai cạnh dài để được hai hình vuông. Chu vi mỗi hình vuông là bao nhiêu xăng-ti-mét?
  \shortans{16}
  \loigiai{
    Mỗi hình vuông có cạnh 4 cm. Chu vi là $4 \times 4 = 16\text{ cm}$.
  }
\end{ex}

% Câu 4
\begin{ex}
  Người ta ghép $6$ hình tam giác đều cạnh $5\text{ cm}$ thành một hình lục giác đều. Chu vi hình lục giác đều đó là bao nhiêu xăng-ti-mét?
  \shortans{30}
  \loigiai{
    Chu vi là $5 \times 6 = 30\text{ cm}$.
  }
\end{ex}

% Câu 5
\begin{ex}
  Dùng $4$ viên gạch hình vuông cạnh $20\text{ cm}$ để ghép thành một hình vuông lớn hơn. Diện tích hình vuông lớn là bao nhiêu xăng-ti-mét vuông?
  \shortans{1600}
  \loigiai{
    Cạnh hình vuông lớn là $40\text{ cm}$. Diện tích là $40 \times 40 = 1\,600\text{ cm}^2$.
  }
\end{ex}

% Câu 6
\begin{ex}
  Cắt một hình vuông có diện tích $64\text{ cm}^2$ thành $4$ hình tam giác bằng nhau. Diện tích mỗi hình tam giác là bao nhiêu xăng-ti-mét vuông?
  \shortans{16}
  \loigiai{
    $64 : 4 = 16\text{ cm}^2$.
  }
\end{ex}

% Câu 7
\begin{ex}
  Một tấm bìa hình vuông cạnh $10\text{ cm}$. Người ta cắt bốn góc bốn hình vuông nhỏ cạnh $2\text{ cm}$. Diện tích phần bìa còn lại là bao nhiêu xăng-ti-mét vuông?
  \shortans{84}
  \loigiai{
    Diện tích hình vuông lớn là 100. Bốn hình vuông nhỏ là $4 \times 4 = 16$. Còn lại $100 - 16 = 84\text{ cm}^2$.
  }
\end{ex}

% Câu 8
\begin{ex}
  Người ta dùng các que tính dài $5\text{ cm}$ để xếp thành một hình thoi có cạnh là 2 que tính. Chu vi hình thoi đó là bao nhiêu xăng-ti-mét?
  \shortans{40}
  \loigiai{
    Cạnh hình thoi là $5 \times 2 = 10\text{ cm}$. Chu vi là $10 \times 4 = 40\text{ cm}$.
  }
\end{ex}

% Câu 9
\begin{ex}
  Dùng một sợi dây thép dài $48\text{ cm}$ để uốn thành một hình vuông. Diện tích hình vuông tạo thành là bao nhiêu xăng-ti-mét vuông?
  \shortans{144}
  \loigiai{
    Cạnh hình vuông là $48 : 4 = 12\text{ cm}$. Diện tích là $12 \times 12 = 144\text{ cm}^2$.
  }
\end{ex}

% Câu 10
\begin{ex}
  Dùng một sợi dây thép dài $48\text{ cm}$ để uốn thành một hình tam giác đều. Độ dài mỗi cạnh của tam giác đều là bao nhiêu xăng-ti-mét?
  \shortans{16}
  \loigiai{
    Cạnh là $48 : 3 = 16\text{ cm}$.
  }
\end{ex}

% Câu 11
\begin{ex}
  Một hình chữ nhật được ghép từ 3 hình vuông bằng nhau, chu vi của mỗi hình vuông là $20\text{ cm}$. Diện tích hình chữ nhật là bao nhiêu xăng-ti-mét vuông?
  \shortans{75}
  \loigiai{
    Cạnh hình vuông là $20 : 4 = 5\text{ cm}$. Diện tích 1 hình vuông là 25. 3 hình vuông là $25 \times 3 = 75\text{ cm}^2$.
  }
\end{ex}

% Câu 12
\begin{ex}
  Một mảnh giấy hình tròn bán kính $10\text{ cm}$ được gập đôi hai lần tạo thành 4 phần bằng nhau. Diện tích mỗi phần là bao nhiêu xăng-ti-mét vuông? (Lấy $\pi = 3{,}14$, làm tròn đến hàng đơn vị).
  \shortans{79}
  \loigiai{
    Diện tích hình tròn là 314. Mỗi phần là $314 : 4 = 78{,}5\text{ cm}^2$. Làm tròn đến hàng đơn vị được 79.
  }
\end{ex}

% Câu 13
\begin{ex}
  Cắt một hình chữ nhật có kích thước $10\text{ cm} \times 6\text{ cm}$ thành các hình vuông nhỏ cạnh $2\text{ cm}$. Cắt được tất cả bao nhiêu hình vuông?
  \shortans{15}
  \loigiai{
    Diện tích lớn là 60. Diện tích hình vuông là 4. Số hình là $60 : 4 = 15$.
  }
\end{ex}

% Câu 14
\begin{ex}
  Cần ít nhất bao nhiêu hình tam giác đều bằng nhau để ghép thành một hình tam giác đều lớn hơn?
  \shortans{4}
  \loigiai{
    Cần ít nhất 4 hình tam giác đều nhỏ.
  }
\end{ex}

% Câu 15
\begin{ex}
  Cần ít nhất bao nhiêu hình vuông bằng nhau để ghép thành một hình vuông lớn hơn?
  \shortans{4}
  \loigiai{
    Cần ít nhất 4 hình vuông nhỏ.
  }
\end{ex}

% Câu 16
\begin{ex}
  Một hình chữ nhật có chiều dài gấp 3 lần chiều rộng và được chia thành 3 hình vuông bằng nhau có tổng diện tích là $108\text{ cm}^2$. Cạnh của mỗi hình vuông là bao nhiêu xăng-ti-mét?
  \shortans{6}
  \loigiai{
    Diện tích mỗi hình vuông là $108 : 3 = 36\text{ cm}^2$. Cạnh là $6\text{ cm}$.
  }
\end{ex}

% Câu 17
\begin{ex}
  Ở câu trên, chu vi của hình chữ nhật ban đầu là bao nhiêu xăng-ti-mét?
  \shortans{48}
  \loigiai{
    Chiều rộng là 6 cm, chiều dài là 18 cm. Chu vi là $(18 + 6) \times 2 = 48\text{ cm}$.
  }
\end{ex}

% Câu 18
\begin{ex}
  Người ta xếp $8$ khối lập phương nhỏ cạnh $1\text{ cm}$ thành một khối lập phương lớn. Thể tích khối lập phương lớn là bao nhiêu xăng-ti-mét khối?
  \shortans{8}
  \loigiai{
    Thể tích là $8\text{ cm}^3$.
  }
\end{ex}

% Câu 19
\begin{ex}
  Ở câu trên, cạnh của khối lập phương lớn là bao nhiêu xăng-ti-mét?
  \shortans{2}
  \loigiai{
    Cạnh khối lập phương lớn là 2 cm.
  }
\end{ex}

% Câu 20
\begin{ex}
  Cần bao nhiêu que tính bằng nhau để xếp thành một hình hộp chữ nhật (khung hình)?
  \shortans{12}
  \loigiai{
    Hình hộp chữ nhật có 12 cạnh, cần 12 que tính.
  }
\end{ex}

% Câu 21
\begin{ex}
  Một ngôi sao 5 cánh được ghép từ 5 hình tam giác cân bằng nhau xung quanh một ngũ giác đều. Biết diện tích mỗi tam giác cân là $8\text{ cm}^2$ và diện tích ngũ giác đều là $20\text{ cm}^2$. Tổng diện tích ngôi sao là bao nhiêu xăng-ti-mét vuông?
  \shortans{60}
  \loigiai{
    Tổng diện tích là $20 + 5 \times 8 = 20 + 40 = 60\text{ cm}^2$.
  }
\end{ex}

% Câu 22
\begin{ex}
  Một tờ giấy hình vuông cạnh $20\text{ cm}$. Người ta cắt một hình tròn có bán kính lớn nhất từ tờ giấy này. Bán kính của hình tròn cắt được là bao nhiêu xăng-ti-mét?
  \shortans{10}
  \loigiai{
    Bán kính lớn nhất bằng nửa cạnh hình vuông: $20 : 2 = 10\text{ cm}$.
  }
\end{ex}

% Câu 23
\begin{ex}
  Ở câu trên, diện tích phần giấy bị cắt bỏ đi là bao nhiêu xăng-ti-mét vuông? (Lấy $\pi = 3{,}14$).
  \shortans{86}
  \loigiai{
    Diện tích hình vuông là $20 \times 20 = 400\text{ cm}^2$. Diện tích hình tròn là $314\text{ cm}^2$. Phần bỏ đi là $400 - 314 = 86\text{ cm}^2$.
  }
\end{ex}

% Câu 24
\begin{ex}
  Vẽ một đường tròn có chu vi $18{,}84\text{ cm}$. Bán kính của đường tròn đó là bao nhiêu xăng-ti-mét? (Lấy $\pi = 3{,}14$).
  \shortans{3}
  \loigiai{
    Bán kính là $18{,}84 : (2 \times 3{,}14) = 3\text{ cm}$.
  }
\end{ex}

% Câu 25
\begin{ex}
  Một khung dây hình tam giác đều có cạnh $8\text{ cm}$. Uốn khung dây này thành một hình chữ nhật có chiều rộng $3\text{ cm}$. Chiều dài hình chữ nhật đó là bao nhiêu xăng-ti-mét?
  \shortans{9}
  \loigiai{
    Chu vi là $8 \times 3 = 24\text{ cm}$. Nửa chu vi là 12 cm. Chiều dài là $12 - 3 = 9\text{ cm}$.
  }
\end{ex}
